

```latex
\documentclass{article}

\usepackage[utf8]{inputenc}
\usepackage[T1]{fontenc}
\usepackage{amsmath,amssymb,amsthm}
\usepackage{hyperref}

\newtheorem{theorem}{Theorem}
\newtheorem{lemma}[theorem]{Lemma}
\newtheorem{proposition}[theorem]{Proposition}
\newtheorem{corollary}[theorem]{Corollary}
\theoremstyle{definition}
\newtheorem{definition}[theorem]{Definition}
\newtheorem{remark}[theorem]{Remark}

\title{Dehn surgery on L-space knots has at most two prime summands}
\author{}
\date{}

\begin{document}

\maketitle

\begin{abstract}
We prove that if \(K\subset S^{3}\) is an L-space knot, then every closed
3-manifold obtained by Dehn surgery on \(K\) has at most two prime summands.
The proof combines the known structure theorem for reducible surgeries on
knots in \(S^{3}\), the classification of possible reducing slopes for
non-cable L-space knots, and the Heegaard Floer obstruction to an integer
homology sphere summand in the exceptional slope \(2g(K)-1\).
\end{abstract}

\section{Conventions and statement}

All knots are oriented knots in \(S^{3}\), and \(S^{3}_{r}(K)\) denotes
Dehn surgery with rational slope \(r\in \mathbb{Q}\cup\{\infty\}\), measured
with respect to the Seifert longitude.  A prime summand means a non-\(S^{3}\)
summand in the prime decomposition.  Orientation reversal does not change the
number of prime summands, so we freely ignore signs of lens spaces when only
the number of summands is relevant.

\begin{theorem}\label{thm:main}
Let \(K\subset S^{3}\) be an L-space knot.  Then for every slope \(r\),
the 3-manifold \(S^{3}_{r}(K)\) has at most two prime summands.
\end{theorem}

\section{Known facts used in the proof}

We collect the external results used below.  We cite the sources by domain.

\begin{lemma}[Gordon--Luecke; Scharlemann; Howie--Sayari--Valdez S\'anchez]
\label{lem:general-reducible}
Let \(K\subset S^{3}\) be a nontrivial knot and suppose
\(S^{3}_{r}(K)\) is reducible.  Then:
\begin{enumerate}
  \item The reducing slope is unique and is an integer \(p\) with
  \(|p|>1\).  Moreover, \(S^{3}_{p}(K)\) contains a nontrivial lens space
  summand.  See
  \href{https://exa.ai/library/publication/zl1m8jt741b}{exa.ai}
  and
  \href{https://ar5iv.labs.arxiv.org/html/0908.2697}{ar5iv.labs.arxiv.org}.
  \item The manifold \(S^{3}_{p}(K)\) has at most three prime summands.
  If it has exactly three prime summands, then it is of the form
  \[
    S^{3}_{p}(K)\cong L(p_{1},a_{1})\# L(p_{2},a_{2})\# Y,
  \]
  where \(p_{1},p_{2}>1\) are coprime, \(|p|=p_{1}p_{2}\), and \(Y\) is an
  integer homology sphere.  See
  \href{https://exa.ai/library/publication/zl1m8jt741b}{exa.ai}
  and
  \href{https://ar5iv.labs.arxiv.org/html/0908.2697}{ar5iv.labs.arxiv.org}.
\end{enumerate}
\end{lemma}

For general background on the question of when knot surgery is non-prime,
see also
\href{https://www.cambridge.org/core/journals/mathematical-proceedings-of-the-cambridge-philosophical-society/article/abs/knot-surgery-and-primeness/29D04C8E36EF794942FA2F0925BE1775}{cambridge.org}.

\begin{lemma}[Reducing slopes of non-cable L-space knots]
\label{lem:HLZ}
Let \(K\) be an L-space knot of genus \(g(K)\) which is not a cable knot.
If \(S^{3}_{r}(K)\) is reducible, then
\[
  r=2g(K)-1.
\]
More precisely, the combined results of Matignon--Sayari and
Hedden--Levine--Zemke imply that a reducing slope of an L-space knot is
either a cabling slope or \(2g(K)-1\); see
\href{https://exa.ai/library/publication/zl1m8jt741b}{exa.ai}
for the Matignon--Sayari cabling alternative and
\href{https://ar5iv.labs.arxiv.org/html/2209.01672}{ar5iv.labs.arxiv.org}
for the statement that, apart from cabling slopes, the only possible reducing
slope for an L-space knot is \(2g(K)-1\).
\end{lemma}

\begin{lemma}[No integer homology sphere summand at slope \(2g-1\)]
\label{lem:no-HS}
Let \(K\) be an L-space knot of genus \(g\).  If
\(S^{3}_{2g-1}(K)\) is reducible, then it cannot have an integer homology
sphere as a connected summand.  Equivalently, if
\[
  S^{3}_{2g-1}(K)\cong Y_{P}\# Y_{Q}
\]
with \(|H_{1}(Y_{P})|=P\), \(|H_{1}(Y_{Q})|=Q\), and \(\gcd(P,Q)=1\), then
\(P,Q>1\).  This is Proposition~4.1 in
\href{https://ar5iv.labs.arxiv.org/html/2209.01672}{ar5iv.labs.arxiv.org}.
\end{lemma}

\begin{lemma}[Reducible surgeries on cable knots]
\label{lem:cable}
Let \(K\subset S^{3}\) be a nontrivial cable knot.  If
\(S^{3}_{r}(K)\) is reducible, then it has at most two prime summands.
In fact, every reducing surgery on a cable knot is the cabling surgery, and
the result is a lens space connected sum an irreducible rational homology
sphere.  See
\href{https://ar5iv.labs.arxiv.org/html/1507.04562}{ar5iv.labs.arxiv.org};
the cabling surgery formula is also recorded in
\href{https://exa.ai/library/publication/zl1m8jt741b}{exa.ai}.
\end{lemma}

\begin{proof}[Proof of Lemma~\ref{lem:cable}]
Let \(K\) be a proper cable of a companion knot \(J\), with cabling slope
\(c\in\mathbb{Z}\).  The standard cabling surgery formula says that
\[
  S^{3}_{c}(K)\cong S^{3}_{q/p}(J)\# L(p,q),
\]
for coprime integers \(p,q\) with \(|p|>1\); see
\href{https://exa.ai/library/publication/zl1m8jt741b}{exa.ai}.
The slope \(q/p\) is non-integral.  By Lemma~\ref{lem:general-reducible}(1),
a reducible surgery on a knot in \(S^{3}\) must have integral slope, so
\(S^{3}_{q/p}(J)\) is irreducible.  The lens space \(L(p,q)\) is prime.
Thus the cabling surgery produces at most two prime summands.

Since the cabling surgery is genuinely reducible for a nontrivial cable knot,
the uniqueness of the reducing slope
(Lemma~\ref{lem:general-reducible}(1)) implies that every reducing surgery on
\(K\) is the cabling surgery.  Hence every reducible surgery on a nontrivial
cable knot has at most two prime summands.
\end{proof}

\section{Proof of Theorem~\ref{thm:main}}

Let \(K\subset S^{3}\) be an L-space knot and let \(M=S^{3}_{r}(K)\).

If \(M\) is irreducible, then it has one prime summand, so there is nothing to
prove.  Assume therefore that \(M\) is reducible.

\subsection*{The unknot}
If \(K\) is the unknot, then \(S^{3}_{r}(K)\) is a lens space for
\(r\neq 0,\infty\), while \(S^{3}_{0}(K)\cong S^{1}\times S^{2}\).  These
manifolds are prime, so the theorem holds.  Henceforth assume \(K\) is
nontrivial.

\subsection*{The cable case}
If \(K\) is a cable knot, then Lemma~\ref{lem:cable} directly implies that
\(M=S^{3}_{r}(K)\) has at most two prime summands.  Thus we may assume that
\(K\) is not a cable knot.

\subsection*{The non-cable case}
Assume now that \(K\) is a non-cable L-space knot of genus \(g=g(K)\).
By Lemma~\ref{lem:HLZ}, the reducing slope is
\[
  r=2g-1.
\]
Thus
\[
  M=S^{3}_{2g-1}(K).
\]

Suppose, for contradiction, that \(M\) has at least three prime summands.
By Lemma~\ref{lem:general-reducible}(2), any reducible surgery on a knot in
\(S^{3}\) has at most three prime summands, and if it has three then it is of
the form
\[
  M\cong L(p_{1},a_{1})\# L(p_{2},a_{2})\# Y,
\]
where \(p_{1},p_{2}>1\) are coprime and \(Y\) is an integer homology sphere.
In particular, \(M\) has an integer homology sphere summand.

But Lemma~\ref{lem:no-HS} says that for an L-space knot,
\(S^{3}_{2g-1}(K)\) cannot have an integer homology sphere summand.  This is a
contradiction.

Therefore \(M\) cannot have three prime summands.  Since it was assumed
reducible, it has exactly two prime summands.  This completes the proof in
the non-cable case.

Combining the unknot case, the cable case, and the non-cable case proves
Theorem~\ref{thm:main}.

\begin{remark}
The argument does not assert that non-cable L-space knots have no reducible
surgeries.  It shows only that if such a reducible surgery exists, it must
occur at slope \(2g(K)-1\) and cannot have three prime summands, because a
three-summand reducible surgery would necessarily contain an integer homology
sphere summand, which is obstructed by Heegaard Floer correction terms as in
Lemma~\ref{lem:no-HS}.
\end{remark}

\begin{thebibliography}{99}

\bibitem{exa}
\href{https://exa.ai/library/publication/zl1m8jt741b}{exa.ai},
\emph{L-space surgeries, genus bounds, and the cabling conjecture}.

\bibitem{0908}
\href{https://ar5iv.labs.arxiv.org/html/0908.2697}{ar5iv.labs.arxiv.org},
\emph{Can Dehn surgery yield three connected summands?}

\bibitem{cambridge}
\href{https://www.cambridge.org/core/journals/mathematical-proceedings-of-the-cambridge-philosophical-society/article/abs/knot-surgery-and-primeness/29D04C8E36EF794942FA2F0925BE1775}{cambridge.org},
\emph{Knot surgery and primeness}.

\bibitem{2209}
\href{https://ar5iv.labs.arxiv.org/html/2209.01672}{ar5iv.labs.arxiv.org},
\emph{Obstructing Reducible Surgeries: Slice Genus and Thickness Bounds}.

\bibitem{1507}
\href{https://ar5iv.labs.arxiv.org/html/1507.04562}{ar5iv.labs.arxiv.org},
\emph{A note on cabled slice knots and reducible surgeries}.

\end{thebibliography}

% Markdown-format source links requested by the prompt:
% [exa.ai](https://exa.ai/library/publication/zl1m8jt741b)
% [ar5iv.labs.arxiv.org](https://ar5iv.labs.arxiv.org/html/0908.2697)
% [cambridge.org](https://www.cambridge.org/core/journals/mathematical-proceedings-of-the-cambridge-philosophical-society/article/abs/knot-surgery-and-primeness/29D04C8E36EF794942FA2F0925BE1775)
% [ar5iv.labs.arxiv.org](https://ar5iv.labs.arxiv.org/html/2209.01672)
% [ar5iv.labs.arxiv.org](https://ar5iv.labs.arxiv.org/html/1507.04562)

\end{document}
```